\documentclass[12pt]{article}
\usepackage{fullpage}
\usepackage[T1]{fontenc}
\usepackage[utf8]{inputenc}
\usepackage{lmodern}
\usepackage{microtype}
\usepackage{amsmath,amssymb,amsthm}
\usepackage{hyperref}

\newtheorem{theorem}{Theorem}
\newtheorem{lemma}{Lemma}
\newtheorem{proposition}{Proposition}
\newtheorem{corollary}{Corollary}
\theoremstyle{definition}
\newtheorem{remark}{Remark}

\DeclareMathOperator{\tr}{tr}
\DeclareMathOperator{\rank}{rank}
\newcommand{\bR}{\mathbb{R}}
\newcommand{\E}{\mathbb{E}}
\newcommand{\eps}{\varepsilon}
\newcommand{\psd}{\succeq 0}

\title{Problem 6: $\eps$-Light Subsets in Graphs\\
\large Partial solution with an honest account of the remaining gap}
\date{}

\begin{document}
\maketitle

\paragraph{Statement.}
For a graph $G=(V,E)$ on $n=|V|$ vertices let $L$ be its Laplacian and, for $S\subseteq V$, let
$L_S$ be the Laplacian of $G_S=(V,E(S,S))$, the subgraph keeping only edges with both endpoints in
$S$. Call $S$ \emph{$\eps$-light} if $\eps L-L_S\psd$. The question is whether there is a universal
constant $c>0$ such that for every graph $G$ and every $\eps\in(0,1)$, $V$ contains an $\eps$-light
subset $S$ with $|S|\ge c\,\eps\,n$.

\paragraph{Summary of what is proved here.}
\begin{itemize}
\item[(A)] The conjectured answer is \textbf{Yes}, and we reduce it cleanly to a single statement
about partitioning positive semidefinite matrices (Section~\ref{sec:reduce}).
\item[(B)] We prove rigorously the structural backbone: the algebraic identities, the
\emph{decoupling inequality} (Lemma~\ref{lem:decouple}), which converts the quadratic lightness
condition into a linear sum of positive semidefinite matrices, and the leverage bounds
(Section~\ref{sec:basics}).
\item[(C)] We prove the bound is \emph{order-optimal} by computing the complete graph \emph{exactly}:
in $K_n$ the largest $\eps$-light set has size exactly $\lfloor\eps n\rfloor$, so no constant $c>1$ is
possible (Proposition~\ref{prop:Kn}). We also prove the two free constructions (single vertex;
independent set).
\item[(D)] We give a precise, correct reduction of the full theorem to the matrix form of Weaver's
$\mathrm{KS}_r$ theorem (Marcus--Spielman--Srivastava), and we identify exactly the one step we could
\emph{not} complete: controlling edges of large effective resistance. We are explicit that, because
of this gap, the universal constant $c$ is \emph{not} fully established here
(Remark~\ref{rem:gap}).
\end{itemize}

Throughout, ``$A\psd$'' and ``$A\preceq B$'' denote the Loewner order, and $\lambda_{\max}(A)$ is the
largest eigenvalue of a symmetric matrix $A$.

%======================================================================
\section{Algebraic preliminaries and the decoupling inequality}\label{sec:basics}

For each edge $e=\{u,v\}\in E$ fix an orientation and set $b_e=\mathbf 1_u-\mathbf 1_v\in\bR^V$ (the
matrix $b_eb_e^\top$ is independent of the orientation). Then
\begin{equation}\label{eq:Ldecomp}
L=\sum_{e\in E} b_e b_e^\top,
\qquad
L_S=\sum_{e=\{u,v\}\in E}\mathbf 1[u\in S]\,\mathbf 1[v\in S]\,b_e b_e^\top .
\end{equation}
Both are positive semidefinite; the missing terms in $L_S$ are positive semidefinite, so
\begin{equation}\label{eq:LSleqL}
0\preceq L_S\preceq L\qquad\text{for every }S\subseteq V .
\end{equation}
In particular $\ker L\subseteq\ker L_S$, $S=V$ is $1$-light, and any independent set $S$ has $L_S=0$
and is $\eps$-light for all $\eps\ge0$.

Let $L^{+}$ be the Moore--Penrose pseudoinverse and $\Pi=L^{+/2}LL^{+/2}$ the orthogonal projection
onto $\operatorname{range}(L)$, where $L^{+/2}\psd$ is the square root of $L^{+}$. For an edge $e$,
its \emph{effective resistance} is
\begin{equation}\label{eq:resistance}
r_e:=b_e^\top L^{+}b_e=\big\|L^{+/2}b_e\big\|^2 .
\end{equation}

\begin{lemma}[Standard resistance facts]\label{lem:resfacts}
$0\le r_e\le 1$ for every $e\in E$, and $\sum_{e\in E}r_e=\rank(L)=n-k_0\le n-1$, where $k_0\ge1$ is
the number of connected components of $G$.
\end{lemma}
\begin{proof}
Nonnegativity is clear from $r_e=\|L^{+/2}b_e\|^2$. For $r_e\le1$, write
$L=b_eb_e^\top+R$ with $R=\sum_{f\ne e}b_fb_f^\top\psd$. Since $b_e\in\operatorname{range}(L)$ we have
$LL^{+}b_e=\Pi b_e=b_e$, hence
\[
r_e=b_e^\top L^{+}b_e=(L^{+}b_e)^\top L (L^{+}b_e)
=\big((L^{+}b_e)^\top b_e\big)^2+(L^{+}b_e)^\top R (L^{+}b_e)
=r_e^2+(L^{+}b_e)^\top R (L^{+}b_e),
\]
using $(L^{+}b_e)^\top b_e=r_e$. As $R\psd$, $r_e\ge r_e^2$, so $r_e\le1$. For the sum, by
\eqref{eq:Ldecomp} and cyclicity of the trace,
$\sum_e r_e=\tr\!\big(L^{+}\sum_e b_eb_e^\top\big)=\tr(L^{+}L)=\tr(\Pi)=\rank(L)=n-k_0$.
\end{proof}

For $u\in V$ define the \emph{leverage load} $\tau_u:=\sum_{e:u\in e}r_e$. Counting each edge at its
two endpoints and using Lemma~\ref{lem:resfacts},
\begin{equation}\label{eq:tau}
\sum_{u\in V}\tau_u=2\sum_{e\in E}r_e\le 2(n-1).
\end{equation}

\subsection*{The decoupling inequality}

The lightness condition $L_S\preceq\eps L$ is \emph{quadratic} in the $0/1$ choice of vertices,
since an edge survives in $L_S$ only when both endpoints are chosen. The next lemma replaces $L_S$ by
a quantity that is \emph{linear} in the selection: a sum of positive semidefinite matrices, one per
chosen vertex. This is the key reduction making partition methods applicable.

Fix any assignment $\phi:E\to V$ with $\phi(e)\in e$ for all $e$, and set
\begin{equation}\label{eq:Mu}
M_u:=\sum_{e:\,\phi(e)=u} b_e b_e^\top\ \psd ,\qquad\text{so}\qquad \sum_{u\in V}M_u=L .
\end{equation}

\begin{lemma}[Decoupling]\label{lem:decouple}
For every $S\subseteq V$, \ $\displaystyle L_S\ \preceq\ \sum_{u\in S} M_u .$
\end{lemma}
\begin{proof}
By \eqref{eq:Ldecomp} and \eqref{eq:Mu},
\[
\sum_{u\in S}M_u-L_S
=\sum_{e=\{u,v\}\in E}\Big(\mathbf 1[\phi(e)\in S]-\mathbf 1[u\in S]\,\mathbf 1[v\in S]\Big)b_eb_e^\top .
\]
For each $e=\{u,v\}$, $\phi(e)\in\{u,v\}$, so if both $u,v\in S$ then $\phi(e)\in S$; thus each scalar
coefficient lies in $\{0,1\}$, in particular is $\ge0$. A nonnegative combination of positive
semidefinite matrices is positive semidefinite.
\end{proof}

After congruence by $L^{+/2}$ (which preserves the Loewner order), with
$A_u:=L^{+/2}M_uL^{+/2}\psd$ and $\sum_{u}A_u=L^{+/2}LL^{+/2}=\Pi\preceq I$, Lemma~\ref{lem:decouple}
gives
\begin{equation}\label{eq:decouple-norm}
L^{+/2}L_S L^{+/2}\ \preceq\ \sum_{u\in S}A_u\qquad(S\subseteq V).
\end{equation}

\begin{lemma}\label{lem:Aunorm}
$\|A_u\|\le\tau_u$ for every $u\in V$.
\end{lemma}
\begin{proof}
$A_u=\sum_{e:\phi(e)=u}(L^{+/2}b_e)(L^{+/2}b_e)^\top$ is a sum of rank-one positive semidefinite
matrices, so $\|A_u\|\le\sum_{e:\phi(e)=u}\|L^{+/2}b_e\|^2=\sum_{e:\phi(e)=u}r_e\le\sum_{e:u\in e}
r_e=\tau_u$.
\end{proof}

Finally we translate lightness into a spectral inequality.

\begin{lemma}\label{lem:equiv}
For $S\subseteq V$: $S$ is $\eps$-light $\iff$ $\lambda_{\max}\!\big(L^{+/2}L_SL^{+/2}\big)\le\eps$.
\end{lemma}
\begin{proof}
Let $x\in\bR^V$, $x=x_0+x_1$ with $x_0\in\ker L$, $x_1\in\operatorname{range}(L)$. By
\eqref{eq:LSleqL}, $\ker L\subseteq\ker L_S$, so $L_S x_0=0=Lx_0$; hence $x^\top L_S x=x_1^\top L_S x_1$
and $x^\top L x=x_1^\top L x_1$. Write $x_1=L^{+/2}y$ with $y\in\operatorname{range}(L)$ (possible since
$L^{+/2}$ is a bijection on $\operatorname{range}(L)$). Then $x^\top L_S x=y^\top(L^{+/2}L_SL^{+/2})y$
and $x^\top Lx=y^\top\Pi y=\|y\|^2$. Thus $\eps L-L_S\psd$ on $\bR^V$ iff
$y^\top(L^{+/2}L_SL^{+/2})y\le\eps\|y\|^2$ for all $y\in\operatorname{range}(L)$; since
$L^{+/2}L_SL^{+/2}$ vanishes on $\ker L=\operatorname{range}(L)^\perp$, this is equivalent to
$\lambda_{\max}(L^{+/2}L_SL^{+/2})\le\eps$.
\end{proof}

%======================================================================
\section{Two free constructions and the exact complete-graph bound}\label{sec:tight}

\begin{proposition}[Free lower bounds]\label{prop:free}
For every graph $G$ and every $\eps\ge0$:
(a) every singleton $S=\{v\}$ is $\eps$-light; (b) every independent set $S$ is $\eps$-light; in
particular there is an $\eps$-light set of size $\alpha(G)$, the independence number.
\end{proposition}
\begin{proof}
In both cases $E(S,S)=\varnothing$, so $L_S=0$ and $\eps L-L_S=\eps L\psd$.
\end{proof}

These already settle two regimes: if $n\le 1/\eps$ then a singleton has size $1\ge\eps n$; and if
$\alpha(G)\ge\eps n$ then an independent set works. The difficult case is dense graphs (small
$\alpha$) with $n$ large.

\begin{proposition}[Complete graph, exact]\label{prop:Kn}
In $K_n$ a set $S$ with $|S|=k$ is $\eps$-light if and only if $k\le\eps n$. Hence the largest
$\eps$-light set in $K_n$ has size $\max\{1,\lfloor\eps n\rfloor\}$. In particular the constant $c$ in
the problem cannot exceed $1$, and the order $\Theta(\eps n)$ is correct.
\end{proposition}
\begin{proof}
For $K_n$, $L=nI-J$ with $J$ the all-ones matrix. Take $S$ to be the first $k$ coordinates. The
induced subgraph $G_S$ is $K_k$ on $S$, so $L_S=kP_S-J_S$, where $P_S$ is the diagonal $0/1$
projection onto coordinates of $S$ and $J_S$ is the all-ones matrix on $S\times S$ (zero elsewhere).
Thus
\[
M:=\eps L-L_S=\eps(nI-J)-(kP_S-J_S).
\]
We diagonalise $M$ exactly. Write $a:=\mathbf 1_S$ and $b:=\mathbf 1_{S^c}$ (the indicator vectors),
so $\mathbf 1=a+b$, $\|a\|^2=k$, $\|b\|^2=n-k$, $a^\top b=0$, $J=(a+b)(a+b)^\top$, $J_S=aa^\top$, and
$P_S a=a$, $P_S b=0$.

\emph{Eigenspace 1.} If $x\in\bR^S$ with $a^\top x=0$ (so $x$ is supported on $S$ and orthogonal to
$a$), then $Jx=(a+b)(a+b)^\top x=(a^\top x)(a+b)=0$, $P_S x=x$, $J_S x=a(a^\top x)=0$, hence
$Mx=\eps n x-0-kx+0=(\eps n-k)x$. This subspace has dimension $k-1$, giving the eigenvalue
$\eps n-k$ with multiplicity $k-1$.

\emph{Eigenspace 2.} If $x\in\bR^{S^c}$ with $b^\top x=0$, then $Jx=0$, $P_Sx=0$, $J_Sx=0$, so
$Mx=\eps n x$. This subspace has dimension $(n-k)-1$, giving eigenvalue $\eps n$ with multiplicity
$n-k-1$.

\emph{The $2$-dimensional remainder.} The orthogonal complement of the two eigenspaces above is
$\operatorname{span}\{a,b\}$, which is $M$-invariant: indeed,
\[
Ma=\eps n a-\eps(a+b)(a^\top(a+b))-ka+a(a^\top a)
=\eps n a-\eps k(a+b)-ka+ka=(\eps n-\eps k)a-\eps k\,b,
\]
\[
Mb=\eps n b-\eps(a+b)(b^\top(a+b))-0+0
=\eps n b-\eps(n-k)(a+b)=-\eps(n-k)a+\eps k\,b .
\]
In the basis $\{a,b\}$ the operator $M|_{\operatorname{span}\{a,b\}}$ has matrix
\[
\begin{pmatrix}\eps(n-k) & -\eps(n-k)\\[2pt] -\eps k & \eps k\end{pmatrix},
\]
whose characteristic polynomial is $\lambda^2-\eps n\,\lambda=\lambda(\lambda-\eps n)$ (trace
$\eps(n-k)+\eps k=\eps n$, determinant $\eps^2 k(n-k)-\eps^2 k(n-k)=0$). Hence the two remaining
eigenvalues are $0$ and $\eps n$.

Collecting all $n$ eigenvalues, the spectrum of $M=\eps L-L_S$ is
\[
\{\,0\,\}\ \cup\ \{\,\eps n\ \text{with multiplicity}\ n-k\,\}\ \cup\ \{\,\eps n-k\ \text{with
multiplicity}\ k-1\,\}.
\]
All eigenvalues except possibly $\eps n-k$ are $\ge0$. Therefore $M\psd$ iff $\eps n-k\ge0$, i.e. iff
$k\le\eps n$. (When $k=1$ there is no eigenvalue $\eps n-k$, consistent with a singleton being light.)
The largest $k$ with $k\le\eps n$ is $\lfloor\eps n\rfloor$ (and a singleton always works), giving the
stated maximum size and the constant claims.
\end{proof}

\begin{remark}
Proposition~\ref{prop:Kn} shows that if the answer is affirmative the best constant is exactly $c=1$
(attained up to the floor by cliques in $K_n$), and that the order $\eps n$ cannot be improved.
\end{remark}

%======================================================================
\section{Reduction of the full theorem and the remaining gap}\label{sec:reduce}

\subsection*{The target reduction}

Combining Lemma~\ref{lem:equiv} with the decoupling inequality \eqref{eq:decouple-norm}: if we can
find $S\subseteq V$ of size $\ge c\eps n$ with
\begin{equation}\label{eq:goal}
\Big\|\sum_{u\in S}A_u\Big\|\le\eps,
\end{equation}
then $\lambda_{\max}(L^{+/2}L_SL^{+/2})\le\eps$ by \eqref{eq:decouple-norm}, so $S$ is $\eps$-light.
Thus the problem reduces to finding a \emph{large} index set $S$ on which the partial sum of the
positive semidefinite matrices $A_u$ (which satisfy $\sum_u A_u\preceq I$, and individually
$\|A_u\|\le\tau_u$) has spectral norm $\le\eps$.

The natural tool is the matrix form of Weaver's $\mathrm{KS}_r$ theorem.

\begin{theorem}[Matrix $\mathrm{KS}_r$; Marcus--Spielman--Srivastava, \emph{Interlacing Families II},
Ann.\ of Math.\ 182 (2015)]\label{thm:KS}
Let $r\ge2$ and $\delta>0$. Let $B_1,\dots,B_m\psd$ be Hermitian $d\times d$ with $\sum_i B_i\preceq I$
and $\|B_i\|\le\delta$ for all $i$. Then there is a partition $\{1,\dots,m\}=X_1\cup\cdots\cup X_r$
with $\big\|\sum_{i\in X_j}B_i\big\|\le(\tfrac1{\sqrt r}+\sqrt\delta)^2$ for every $j$.
\end{theorem}

To reach the target \eqref{eq:goal} with $r=\Theta(1/\eps)$ blocks (so that some block has size
$\ge|V'|/r=\Omega(\eps n)$ by averaging, after discarding the at most $2(n-1)/T$ vertices with
$\tau_u\ge T$ for a constant $T$, which leaves a constant fraction of $V$ by \eqref{eq:tau}), one needs
the per-matrix norm bound $\delta$ to satisfy $(\tfrac1{\sqrt r}+\sqrt\delta)^2\le\eps$. With
$r\asymp 1/\eps$ the first term is $\Theta(\sqrt\eps)$, so this \emph{requires} $\delta=O(\eps)$: we
must partition into atoms each of spectral norm $O(\eps)$.

The atoms of $A_u=\sum_{e:\phi(e)=u}(L^{+/2}b_e)(L^{+/2}b_e)^\top$ are the rank-one matrices
$(L^{+/2}b_e)(L^{+/2}b_e)^\top$ of norm $r_e$. Theorem~\ref{thm:KS} therefore applies to the
\emph{edge atoms} with $\delta=\max_e r_e$.

\begin{remark}[The remaining gap]\label{rem:gap}
This reduction is rigorous, but it reaches the target only when every atom has norm $O(\eps)$, i.e.
only for edges of effective resistance $r_e=O(\eps)$. Edges with $r_e>\eps$ (\emph{heavy} edges) are a
genuine obstruction: a single atom of norm $r_e$ forces \emph{some} block of \emph{any} partition to
have norm $\ge r_e>\eps$, so no partition into blocks of norm $\le\eps$ exists when heavy edges are
present. To finish, one must dispose of heavy edges separately and recombine. I was not able to
complete this step, and I record honestly that the two natural attempts fail:
\begin{itemize}
\item \emph{Independent set of the heavy subgraph.} The subgraph $H$ formed by heavy edges need not
have a large independent set; explicit dense examples show $\alpha(H)$ can be a small constant fraction
of $n$ (checked numerically), so the heavy part cannot be routed through Proposition~\ref{prop:free}.
The simple count $|E(H)|\le(n-1)/\eps$ (from $\sum_e r_e\le n-1$) is too weak, since it may exceed $n$.
\item \emph{Matrix concentration.} Independent vertex sampling with Tropp's matrix Chernoff inequality
does \emph{not} help: with per-term norm $\|A_u\|\le T=O(1)$ and target $\eps<1$, the effective number
of samples $\mu/T<1$ lies far below $\log d$ ($d=n$), so the dimensional union bound cannot be beaten
unless the sampling rate is exponentially small in $1/\eps$, which destroys linear size. (In
particular there is \emph{no} valid bound of the shape $\eps n/\log n$ from this route; a claim to that
effect would be incorrect.)
\end{itemize}
Resolving the heavy edges appears to require applying the interlacing-families method directly to the
original \emph{vertex} selection (a quadratic, higher-degree variant of $\mathrm{KS}_r$), which I could
not carry out correctly here.
\end{remark}

\subsection*{Numerical evidence for the affirmative answer}

Direct computation supports the affirmative answer with an absolute constant: for $K_n$ the optimal
$\eps$-light set has size exactly $\lfloor\eps n\rfloor$ (Proposition~\ref{prop:Kn}); for dense
Erd\H{o}s--R\'enyi graphs $G(n,1/2)$ and random regular graphs, randomly sampled vertex sets of rate
$\Theta(\eps)$ are $\eps$-light with sizes $\Theta(\eps n)$ and ratios that remain bounded away from
$0$ as $n$ grows (verified for $n$ up to a few hundred). No instance was found that drives the
optimal ratio (size)$/(\eps n)$ to $0$.

%======================================================================
\section{Conclusion}\label{sec:concl}

We have rigorously established: the algebraic framework and the decoupling inequality
(Section~\ref{sec:basics}), reducing $\eps$-lightness of $S$ to the linear spectral bound
\eqref{eq:goal}; the free constructions and the \emph{exact} extremal answer $\lfloor\eps n\rfloor$
for the complete graph (Section~\ref{sec:tight}), which fixes the optimal constant at $c=1$ and
confirms the order $\Theta(\eps n)$; and a correct reduction of the general theorem to the matrix
$\mathrm{KS}_r$ theorem, isolating the single missing ingredient: control of edges of large effective
resistance (Remark~\ref{rem:gap}).

The expected answer to the problem is \textbf{Yes}. We are explicit, however, that a complete proof of
the universal constant $c$ is \emph{not} achieved here, because of the gap in Remark~\ref{rem:gap}.
\end{document}
